% Chapter 1: Primitives
% Abridged from book ch01_objects.tex, keeping only the vocabulary used by
% the characterization chapters (Ch 2-6). Full hub-node exposition of the
% concept class and the MDL/MML/Bayesian analogies is deferred to the book.

\chapter{Primitives}\label{chap:primitives}

Every learning problem begins with the same question: given data drawn
from an unknown source, find a rule that predicts well on future data.
This chapter introduces the mathematical vocabulary for making that
question precise. The definitions are unremarkable individually.  The
mathematical content begins when we ask how they compose.

\section{Atomic vocabulary}\label{sec:atoms}

\begin{defn}[Domain, Label, Concept]
\label{def:domain-label-concept}
  \lean{Concept}
  \leanok
A \emph{domain} $X$ is a set whose elements are called \emph{instances}.
A \emph{label set} $Y$ is a set of possible outputs; in binary
classification, $Y = \{0, 1\}$. A \emph{concept} is a function
$c : X \to Y$.

Equivalently, when $Y = \{0,1\}$, a concept $c$ can be identified with
the subset $\{x \in X : c(x) = 1\} \subseteq X$. Both perspectives are
used throughout the literature. We adopt the function view as primary
and the set view when it simplifies combinatorial arguments (as in
shattering, \Cref{chap:pac}).
\end{defn}

\begin{defn}[Hypothesis, Target Concept, Proper vs.\ Improper]
\label{def:hypothesis-target}
A \emph{hypothesis} $h : X \to Y$ is a candidate prediction rule that a
learning algorithm might output. The \emph{target concept}
$c^* \in \mathcal{C}$ is the specific concept that generated the
training data. Learning is \emph{proper} if the algorithm's output is
constrained to lie in $\mathcal{C}$, and \emph{improper} if it may use
a larger hypothesis space $\mathcal{H} \supseteq \mathcal{C}$.
\end{defn}

\section{The concept class}\label{sec:concept-class}

\begin{defn}[Concept Class]
\label{def:concept-class}
  \lean{ConceptClass}
  \leanok
A \emph{concept class} $\mathcal{C} \subseteq Y^X$ is a collection of
concepts over a common domain. The class is the primary object of
study: every major question in learning theory is a question about
concept classes.
\end{defn}

The concept class is the node every complexity measure in the book
connects to: VC dimension, Littlestone dimension, Rademacher
complexity, covering number, growth function, mind-change dimension.
It is the object whose structure determines what is learnable.

\section{Batch learner}\label{sec:batch-learner}

\begin{defn}[Batch Learner]
\label{def:batch-learner}
  \lean{BatchLearner}
  \leanok
A \emph{batch learner} over $(X, Y)$ is a bundled structure carrying a
hypothesis set $\mathcal{H} \subseteq Y^X$, a learning function
$\mathcal{A}$ that takes a finite training sample
$S : \mathrm{Fin}\,m \to X \times Y$ and returns a hypothesis in
$\mathcal{H}$, and a proof that the output lies in $\mathcal{H}$.
\end{defn}

The kernel formalizes the batch learner as a structure rather than as
a bare function so that the constraint ``output lies in the hypothesis
set'' travels with the algorithm. Chapter~\ref{chap:online} will
introduce a sequential learner with its own bundled type; the contrast
between batch and sequential learning is the subject of
Chapters~\ref{chap:pac} and~\ref{chap:online}.

\section{Sample, loss, and error}\label{sec:sample-loss}

\begin{defn}[I.i.d.\ sample]
\label{def:iid-sample}
For a distribution $D$ on $X$ and a target concept $c$, an
\emph{i.i.d.\ sample} of size $m$ is a tuple
$S = ((x_1, c(x_1)), \dots, (x_m, c(x_m)))$ with each $x_i \sim D$
drawn independently.
\end{defn}

\begin{defn}[Zero-one loss]
\label{def:zero-one-loss}
For $h, c : X \to \{0,1\}$ and $x \in X$, the \emph{zero-one loss} is
$\ell(h, c, x) = \mathbf{1}[h(x) \neq c(x)]$.
\end{defn}

\begin{defn}[Empirical and true error]
\label{def:errors}
For a hypothesis $h$, a target $c$, a distribution $D$, and a sample
$S$ of size $m$:
\begin{align*}
  \emprisk_S(h) &= \frac{1}{m} \sum_{i=1}^{m} \mathbf{1}[h(x_i) \neq c(x_i)]
    && \text{(empirical error)}, \\
  \risk_D(h) &= \Pr_{x \sim D}\!\left[ h(x) \neq c(x) \right]
    && \text{(true error)}.
\end{align*}
\end{defn}
